\begin{center}
\SourceLeaf{027c}\pdfbookmark[2]{10.}{article-010}10.
\end{center}
Aequationes art. 8 omnia continent, quae ad computum anomaliae excentricae et mediae e vera, vel excentricae et verae e media requiruntur. Pro deducenda excentrica e vera vulgo formula VII adhibetur; plerumque tamen praestat ad hunc finem aequ. X uti, praesertim quoties excentricitas non nimis magna est, in quo casu $E$ per X maiori praecisione computari potest, quam per VII. Praeterea adhibita aequatione X, logarithmus sinus $E$, qui in XII requiritur, protinus per aequationem VIII habetur, quem adhibita VII e tabulis arcessere oporteret; si igitur in illa methodo hic logarithmus etiam e tabulis desumitur, simul calculi recte instituti confirmatio hinc obtinetur. Huiusmodi calculi examina et comprobationes magni semper sunt aestimanda, quibus igitur consulere in omnibus methodis in hoc opere tradendis, ubi quidem commode fieri potest, assiduae nobis ubique curae erit. — Ad maiorem illustrationem exemplum complete calculatum adiungimus.

Data sint $v = 310^0\,55^{\prime}\,29^{\prime\prime}{,}64$, $\varphi = 14^0\,12^{\prime}\,1^{\prime\prime}{,}87$, $\log r = 0{,}3307640$; quaeruntur $p$, $a$, $E$ et $M$.

\SourceLeaf{028}
\begin{center}
\begin{tabular}{@{}ll@{}}
\makebox[16em][l]{$\log\sin\varphi$\hfill .\,.\,.\,.\,.\,.\,.\,.\,.\,.} & $9{,}389\,7262$\\[.3em]
\makebox[16em][l]{$\log\cos v$\hfill .\,.\,.\,.\,.\,.\,.\,.\,.\,.} & $9{,}816\,2872$\\[.3em]
\hline
 & $9{,}206\,0134$\\[.3em]
\multicolumn{2}{r}{unde $e\cos v=0{,}1606991$}\\[.3em]
\hline
\makebox[16em][l]{$\log(1+e\cos v)$\hfill .\,.\,.\,.\,.\,.} & $0{,}064\,7197$\\[.3em]
\makebox[16em][l]{$\log r$\hfill .\,.\,.\,.\,.\,.\,.\,.\,.\,.\,.\,.\,.} & $0{,}330\,7640$\\[.3em]
\hline
\makebox[16em][l]{$\log p$\hfill .\,.\,.\,.\,.\,.\,.\,.\,.\,.\,.\,.\,.} & $0{,}395\,4837$\\[.3em]
\makebox[16em][l]{$\log\cos\varphi^2$\hfill .\,.\,.\,.\,.\,.\,.\,.\,.} & $9{,}973\,0448$\\[.3em]
\hline
\makebox[16em][l]{$\log a$\hfill .\,.\,.\,.\,.\,.\,.\,.\,.\,.\,.\,.\,.} & $0{,}422\,4389$\\[.3em]
\hline
\makebox[16em][l]{$\log\sin v$\hfill .\,.\,.\,.\,.\,.\,.\,.\,.\,.} & $9{,}878\,2740_n$\footnotemark\\[.3em]
\makebox[16em][l]{$\log\BareRoot{\frac{p}{r}}$\hfill .\,.\,.\,.\,.\,.\,.\,.\,.\,.} & $0{,}032\,3598.5$\\[.3em]
\hline
 & $9{,}845\,9141.5_n$\\[.3em]
\makebox[16em][l]{$\log\sin\frac{1}{2}\varphi$\hfill .\,.\,.\,.\,.\,.\,.\,.\,.} & $9{,}092\,0395$\\[.3em]
\hline
\noalign{\vskip2pt}
\makebox[16em][l]{$\log\sin\frac{1}{2}(v-E)$\hfill .\,.\,.\,.\,.} & $8{,}937\,9536.5_n$\\[.3em]
\end{tabular}
\end{center}
\footnotetext{Litera $n$ logarithmo affixa indicat, numerum cui respondet negativum esse.}

\noindent hinc
\[
\begin{gathered}
\frac{1}{2}(v-E)=-4^0\,58^{\prime}\,22^{\prime\prime}{,}94;\quad v-E=-9^0\,56^{\prime}\,45^{\prime\prime}{,}88;\\
E=320^0\,52^{\prime}\,15^{\prime\prime}{,}52.
\end{gathered}
\]
\noindent Porro fit
\begin{center}
\begin{tabular}{@{}ll@{}}
\makebox[16em][l]{$\log e$\hfill .\,.\,.\,.\,.\,.\,.\,.\,.} & $9{,}389\,7262$\\[.3em]
\makebox[16em][l]{$\log206264{,}8$\hfill .\,.\,.} & $5{,}314\,4251$\\[.3em]
\hline
\makebox[16em][l]{$\log e$ in sec:\hfill .\,.\,.} & $4{,}704\,1513$\\[.3em]
\makebox[16em][l]{$\log\sin E$\hfill .\,.\,.\,.\,.\,.} & $9{,}800\,0767_n$\\[.3em]
\hline
 & $4{,}504\,2280_n$\\[.3em]
\end{tabular}
\end{center}

\begin{center}
Calculus pro $\log\sin E$ per formulam VIII.
\end{center}
\begin{center}
\begin{tabular}{@{}ll@{}}
\makebox[16em][l]{$\log\frac{r}{p}\sin v$\hfill .\,.\,.\,.\,.} & $9{,}813\,5543_n$\\[.3em]
\makebox[16em][l]{$\log\cos\varphi$\hfill .\,.\,.\,.\,.\,.} & $9{,}986\,5224$\\[.3em]
\hline
\makebox[16em][l]{$\log\sin E$\hfill .\,.\,.\,.\,.\,.} & $9{,}800\,0767_n$\\[.3em]
\hline
\end{tabular}
\end{center}

\noindent hinc
\[
\begin{gathered}
e\sin E\ \mathrm{in\ secundis}=-31932^{\prime\prime}{,}14=-8^0\,52^{\prime}\,12^{\prime\prime}{,}14\\
\mathrm{atque}\quad M=329^0\,44^{\prime}\,27^{\prime\prime}{,}66.
\end{gathered}
\]
\noindent Per formulam VII calculus pro $E$ ita se haberet:
\begin{center}
\begin{tabular}{@{}ll@{}}
\multicolumn{2}{l}{$\frac{1}{2}v=155^0\,27^{\prime}\,44^{\prime\prime}{,}82$}\\[.3em]
\makebox[16em][l]{$\log\operatorname{tang}\frac{1}{2}v$\hfill .\,.\,.\,.\,.\,.\,.\,.} & $9{,}659\,4579_n$\\[.3em]
\multicolumn{2}{l}{$45^0-\frac{1}{2}\varphi=37^0\,53^{\prime}\,59^{\prime\prime}{,}065$}\\[.3em]
\makebox[16em][l]{$\log\operatorname{tang}(45^0-\frac{1}{2}\varphi)$\hfill .\,.\,.} & $9{,}891\,2427$\\[.3em]
\hline
\noalign{\vskip2pt}
\makebox[16em][l]{$\log\operatorname{tang}\frac{1}{2}E$\hfill .\,.\,.\,.\,.\,.\,.} & $9{,}550\,7006_n$\\[.3em]
\end{tabular}
\end{center}
unde $\frac{1}{2}E=160^0\,26^{\prime}\,7^{\prime\prime}{,}76$ atque $E=320^0\,52^{\prime}\,15^{\prime\prime}{,}52$ ut supra.
